\chapter{MARCO TEÓRICO Y REFERENCIAL}

\section{ANTECEDENTES DE LA INVESTIGACIÓN}

Los estudios comparativos de frameworks de interfaz de usuario (Banks y
Porcello, 2023; Brown, 2023) documentan diferencias de arquitectura,
rendimiento y ecosistema entre Angular, React y Vue.js en aplicaciones
empresariales, sin abordar las restricciones propias de sectores regulados
como el bancario.

Las investigaciones sobre selección tecnológica en instituciones financieras
(Arner et al., 2017; Chen et al., 2021) evidencian que el cumplimiento
normativo y la gobernanza con múltiples stakeholders condicionan las
decisiones tecnológicas del sector, aunque no proponen instrumentos de
evaluación comparativa de frameworks de interfaz.

Las metodologías genéricas de evaluación de software y de decisión
multicriterio aplicadas a la selección tecnológica (Kumar y Sharma, 2020)
constituyen la base metodológica más cercana a este trabajo; sin embargo,
omiten los requisitos regulatorios bancarios y los criterios emergentes
asociados a la IA generativa, brecha que la presente tesis busca cerrar.

Finalmente, la evidencia reciente sobre generación de código asistida por IA
(Peng et al., 2023; Stack Overflow, 2024) muestra su adopción masiva y su
efecto sobre la productividad y la calidad del código, lo que motiva
incorporar la calidad del código generado y la gobernanza del uso de IA como
criterios de evaluación. En la tesis, estos antecedentes se ampliarán y
sistematizarán con detalle.

\section{BASES TEÓRICAS}

Se enuncian brevemente las bases teóricas sobre las que se construirá la
metodología; su desarrollo extenso corresponderá al marco teórico de la tesis.

\subsection{Tecnologías de interfaz de usuario y frameworks front-end}

La evolución de las interfaces bancarias, desde terminales de texto hasta
experiencias digitales omnicanal (Dapp, 2014; King, 2018), ha consolidado a
Angular, React y Vue.js como las principales alternativas empresariales, con
paradigmas distintos: framework completo, biblioteca de vistas y framework
progresivo, respectivamente. En la tesis se analizarán sus arquitecturas,
modelos de renderizado, ecosistemas y características distintivas en el
contexto bancario.

\subsection{Modelos y criterios de evaluación tecnológica}

Los modelos de calidad de producto software, las técnicas de evaluación
arquitectónica y los enfoques multicriterio ofrecen una base genérica para
comparar tecnologías; no obstante, resultan insuficientes para el sector
financiero al no contemplar sus exigencias regulatorias (Kumar y Sharma,
2020). La tesis adaptará y extenderá estos modelos con criterios técnicos y
regulatorios específicos del sector.

\subsection{Regulaciones y cumplimiento en el sector bancario}

La regulación de servicios financieros digitales, la protección de datos y la
resiliencia operativa (PCI-DSS, DORA y normativas SBS y locales; Arner et al.,
2017) imponen requisitos que condicionan la selección de tecnologías de
interfaz. En la tesis se sistematizarán los requisitos aplicables y su
traducción a criterios evaluables.

\subsection{Procesos decisionales y gobernanza tecnológica}

Las decisiones tecnológicas bancarias involucran comités y stakeholders de
arquitectura, seguridad y cumplimiento (Chen et al., 2021); la fragmentación
de criterios entre estas áreas produce decisiones subóptimas. La tesis
propondrá un proceso estructurado de toma de decisiones con roles y flujos
definidos.

\subsection{Teoría de decisiones multicriterio}

El Proceso Analítico Jerárquico (AHP) descompone un problema decisional en una
jerarquía de criterios y deriva pesos mediante comparaciones pareadas
(Triantaphyllou, 2000); se complementa con métodos de ordenamiento como
TOPSIS y PROMETHEE. La tesis empleará AHP para ponderar los criterios de
evaluación técnicos y regulatorios.

\subsection{Tendencias emergentes e IA generativa}

Los micro-frontends, WebAssembly y la personalización basada en IA configuran
la evolución reciente de las interfaces; a ello se suma la generación de
código asistida por IA, que reduce las diferencias de ecosistema y talento
entre frameworks (Peng et al., 2023). La tesis analizará el efecto moderador
de estas herramientas sobre la ponderación de los criterios de selección.

\section{MARCO CONCEPTUAL}

Se precisan los conceptos básicos de referencia, que serán desarrollados con
amplitud en la tesis.

\textbf{Tecnologías de interfaz de usuario:} conjunto de componentes,
herramientas y estándares con los que una aplicación presenta información y
captura las interacciones del usuario en sus distintos canales.

\textbf{Frameworks de interfaz de usuario:} conjuntos estructurados de
librerías, convenciones y herramientas que estandarizan la construcción de
interfaces; en este estudio, Angular, React y Vue.js.

\textbf{Tecnologías complementarias:} elementos que acompañan al framework en
aplicaciones empresariales: gestión de estado, enrutamiento, renderizado del
lado del servidor, sistemas de diseño y herramientas de prueba.
